\chapter{Galerkin 粗网格算子的一维推导}
\label{app:galerkin}

\begin{keybox}
\textbf{目的：}把\chapref{ch:two-grid}的式~\eqref{eq:galerkin}展开成可逐行检查的一维\textbf{内部网格}计算，验证本书 nodal FD 代表例中的标准 linear prolongation/full-weighting 下，Galerkin coarse operator 的 interior stencil 确实回到标准 $2h$ Poisson 离散。边界行仍取决于 Dirichlet/Neumann/periodic 的离散闭合，不能由这段内部推导自动推出；cell-centered transfer 的几何位置与权重也不同，不由本附录这一套 nodal 公式代表。
\end{keybox}


令 fine-grid operator
\begin{equation}
(A_hu)_j=\frac{-u_{j-1}+2u_j-u_{j+1}}{h^2}.
\end{equation}
取线性 prolongation：
\begin{align}
(Pv)_{2I}&=v_I,\\
(Pv)_{2I+1}&=\frac12(v_I+v_{I+1}).
\end{align}
先计算 $A_hPv$ 在偶点：
\begin{align}
(A_hPv)_{2I}
&=\frac{1}{h^2}
\left[
-v_{I-1/2}^{interp}+2v_I-v_{I+1/2}^{interp}
\right]\\
&=\frac{1}{h^2}
\left[-\frac12(v_{I-1}+v_I)+2v_I-\frac12(v_I+v_{I+1})\right]\\
&=\frac{-v_{I-1}+2v_I-v_{I+1}}{2h^2}.
\end{align}
在奇点 $2I+1$：
\begin{align}
(A_hPv)_{2I+1}
&=\frac{1}{h^2}
\left[-v_I+2\cdot\frac{v_I+v_{I+1}}2-v_{I+1}\right]\\
&=0.
\end{align}

full-weighting restriction 为 $R=\frac12P^T$：
\begin{equation}
(Rw)_I=\frac14w_{2I-1}+\frac12w_{2I}+\frac14w_{2I+1}.
\end{equation}
由于 $w=A_hPv$ 的奇点值为零，
\begin{align}
(RA_hPv)_I
&=\frac12(A_hPv)_{2I}\\
&=\frac{-v_{I-1}+2v_I-v_{I+1}}{4h^2}.
\end{align}
又因为 $H=2h$，
\begin{equation}
4h^2=H^2,
\end{equation}
所以
\begin{equation}
\boxed{A_H=RA_hP
=\frac{1}{H^2}\operatorname{tridiag}(-1,2,-1)}.
\end{equation}
也就是说，对内部三点 stencil，在这个经典组合下 Galerkin 粗算子恰好等于直接在 $2h$ 网格上重新离散得到的内部 Poisson stencil。

这段推导把“Galerkin 是抽象矩阵乘法”和“几何粗网格仍表示同一个低频 PDE 结构”连接起来，但不要把结论外推过头。边界行、变系数、非均匀网格、cell-centered transfer 或 AMG 中，$RA_hP$ 与直接 rediscretization 的逐项一致性都需要单独检查。

\section*{推导检查}
\begin{checkpointbox}
不看正文，重新计算 $(A_hPv)_{2I}$ 与 $(A_hPv)_{2I+1}$。为什么奇点结果为零？这个零值怎样使 $RA_hP$ 恰好得到 $H=2h$ 的三点 Laplacian？
\end{checkpointbox}
